\subsubsection{Resultados esperados [Pendientes: 1. acordar criterios finales de aceptación académica]}\label{sec:transicion-resultados}
La fase deberá cerrar con una versión demostrable corregida, el informe de evaluación E6, el manual y el informe final E7, junto con la matriz de trazabilidad E1 y los artefactos técnicos actualizados que respaldan la entrega. Su cumplimiento se comprobará mediante la reproducción de la demostración, la revisión de resultados frente a los criterios acordados y la identificación explícita de limitaciones o problemas no resueltos. La documentación de una evaluación desfavorable también constituirá evidencia del proyecto y deberá reflejarse en sus conclusiones, en lugar de omitirse para presentar un resultado aparentemente satisfactorio.

La evaluación de comprensión completará principalmente la evidencia de OE4, mientras que la demostración integrada y la trazabilidad final permitirán revisar conjuntamente los aportes de OE1, OE2 y OE3. Esta fase no sustituye las pruebas realizadas durante el desarrollo, sino que consolida su evidencia en una entrega evaluable. La aceptación académica se determinará con el director a partir de los criterios y evidencias consolidados.
